\section{从词嵌入到全模型预训练}

\subsection{分布式假设的两个层次}

课程开篇曾引用 J.R. Firth 的名言：``You shall know a word by the company it keeps.'' 这是\textbf{分布式假设}的经典表述，也是 Word2Vec 的理论基础。

但同一位语言学家还有另一句话：``The complete meaning of a word is always contextual, and no study of meaning apart from a complete context can be taken seriously.'' 这揭示了一个更深层的洞察：词义不仅由其典型的共现伙伴决定，更由其\textbf{具体出现的上下文}决定。

\begin{figure}[H]
\centering
\includegraphics[width=0.95\textwidth]{slides-images/slide-012.jpg}
\caption{从 Word2Vec 到上下文表示的演进：Word2Vec 为 ``record'' 只学习一个向量（混合了名词义和动词义）；上下文表示则能根据具体语境给出不同的向量\protect\footnotemark}
\end{figure}
\footnotetext{Slides 第13页。}

\begin{importantbox}{Word2Vec 的根本局限}
Word2Vec 为每个词学习\textbf{一个固定的向量}，无论它出现在什么上下文中。对于多义词（如英语的 ``record''：动词``记录'' / 名词``唱片''），这个向量只能是所有义项的\textbf{混合体}，无法在特定语境中区分不同含义。
\end{importantbox}

\subsection{旧范式 vs 新范式}

在预训练出现之前，NLP 的典型工作流程是：

\begin{enumerate}
    \item 用 Word2Vec / GloVe 训练词嵌入（预训练的\textbf{仅仅是嵌入层}）
    \item 在嵌入层之上堆叠一个随机初始化的 LSTM 或 Transformer
    \item 用下游任务（如情感分析、机器翻译）的标注数据从头训练上层网络
\end{enumerate}

新范式的革命性变化在于：\textbf{预训练整个网络}（嵌入层 + Transformer），而不仅仅是词嵌入。

\begin{figure}[H]
\centering
\includegraphics[width=0.95\textwidth]{slides-images/slide-013.jpg}
\caption{旧范式 vs 新范式：左侧只预训练词嵌入（绿色），上层网络（灰色）随机初始化；右侧整个网络都经过预训练（全部绿色），然后在下游任务上微调\protect\footnotemark}
\end{figure}
\footnotetext{Slides 第14页。}

这带来了三大优势：

\begin{enumerate}
    \item \textbf{上下文表示}：``record'' 在不同语境中获得不同的向量
    \item \textbf{强参数初始化}：所有参数都从一个``好的起点''开始，而非随机初始化
    \item \textbf{语言概率分布}：模型学到了自然语言的生成概率，可以直接用于文本生成
\end{enumerate}

\subsection{预训练的核心思想：重构输入}

\begin{figure}[H]
\centering
\includegraphics[width=0.95\textwidth]{slides-images/slide-015.jpg}
\caption{预训练的核心思想：``Stanford University is located in \underline{\hspace{1.5cm}}, California.'' ——通过让模型填补空白来学习语言知识\protect\footnotemark}
\end{figure}
\footnotetext{Slides 第16页。}

预训练的核心假设是：如果我们遮住文本的一部分，让神经网络来\textbf{重构被遮住的内容}，那么网络为了做好这个任务，必须学会大量关于语言和世界的知识。

通过这种``填空''任务，模型可以隐式地学到：

\begin{itemize}
    \item \textbf{世界知识}（trivia）：``Stanford University is located in \underline{Stanford/Palo Alto}, California''
    \item \textbf{句法}（syntax）：``I put \underline{the/a} fork down on the table''
    \item \textbf{共指}（coreference）：``The woman walked across the street, checking for traffic over \underline{her} shoulder''
    \item \textbf{语义}（semantics）：``I went to the ocean to see the fish, turtles, seals, and \underline{whales}''
    \item \textbf{情感}（sentiment）：``The movie was \underline{bad/great}''
    \item \textbf{推理}（reasoning）：``Iroh went into the kitchen... Zuko left the \underline{kitchen}''
\end{itemize}

\begin{figure}[H]
\centering
\includegraphics[width=0.95\textwidth]{slides-images/slide-020.jpg}
\caption{推理示例：Iroh 在厨房泡茶，Zuko 站在 Iroh 旁边，因此 Zuko 也在厨房——``Zuko left the \underline{kitchen}'' 需要模型进行空间推理\protect\footnotemark}
\end{figure}
\footnotetext{Slides 第21页。}

\subsection{预训练--微调范式}

\begin{figure}[H]
\centering
\includegraphics[width=0.95\textwidth]{slides-images/slide-023.jpg}
\caption{预训练--微调的两阶段流程：第一阶段在海量无标注文本上做语言建模；第二阶段用少量标注数据在下游任务上微调\protect\footnotemark}
\end{figure}
\footnotetext{Slides 第24页。}

\begin{importantbox}{预训练--微调范式（Pretrain-Finetune Paradigm）}
\begin{enumerate}
    \item \textbf{预训练}：在海量无标注文本上训练，学习语言的通用表示。目标是最小化预训练损失 $\hat{\theta} = \arg\min_\theta \mathcal{L}_{\text{pretrain}}(\theta)$
    \item \textbf{微调}：以 $\hat{\theta}$ 为起点，在下游任务的标注数据上继续训练。梯度下降从 $\hat{\theta}$ 出发，走向对当前任务最优的参数 $\theta^*$
\end{enumerate}
\end{importantbox}

为什么预训练的参数初始化如此重要？Chris Manning 强调：``If you could actually solve this min, the starting point shouldn't matter. But it really, really, really does.'' 直觉上：

\begin{itemize}
    \item 梯度下降在微调过程中会停留在 $\hat{\theta}$ 附近的局部最优
    \item $\hat{\theta}$ 附近的参数空间已经``学会了语言''，因此该区域的局部最优往往\textbf{泛化能力更强}
    \item 无标注数据的量级（万亿词级别）远超标注数据（百万词级别），预训练利用了这个巨大的信息优势
\end{itemize}

\begin{warningbox}{预训练的数据优势不仅在于量}
即使你拥有海量的情感分析标注数据，从头训练的模型也可能不如预训练+微调的模型。原因在于：
\begin{itemize}
    \item 单一任务的标注数据多样性有限，模型容易过拟合到特定分布
    \item 预训练数据涵盖了互联网上极其多样的文本类型，使模型具备更强的\textbf{分布外泛化}（out-of-distribution generalization）能力
    \item 语言建模本身是一个极其困难的任务，要求模型学习语言的方方面面
\end{itemize}
\end{warningbox}

\subsection{本章小结}

从 Word2Vec 的``仅预训练嵌入''到``预训练整个网络''，NLP 经历了一次范式转变。预训练--微调范式的成功基于两个关键因素：（1）海量无标注文本提供了丰富的自监督信号；（2）预训练学到的参数提供了一个泛化能力极强的初始化点。这一范式至今仍是 NLP 的主流方法论。

%% ========== Part 3: Encoder 预训练 ==========

